\hypertarget{appendix-0ec599f8ba27}{}
\section*{Numerical tolerance policies}
\chapter{Numerical tolerance policies}\label{appendix-0ec599f8ba27-numerical-tolerance-policies}

\textbf{Curriculum depth:} Developed · \textbf{First seen:} C004 · \textbf{Latest development:} C005

\section{Definition and mechanism}\label{appendix-0ec599f8ba27-definition-and-mechanism}

Tolerance is a task-chosen allowed discrepancy: absolute in original units, relative to scale, or representable-step based. Machine epsilon is not automatically a norm cutoff. Handle non-finite values and validate tolerances before formulas.

M005 distinguishes absolute tolerance in the quantity's units from dimensionless relative tolerance. Validate policy parameters and classify non-finite values before arithmetic; approximate closeness is not transitive.

\section{Prerequisites}\label{appendix-0ec599f8ba27-prerequisites}

\begin{itemize}
\tightlist
\item
  \hyperref[appendix-390ea641da54]{Machine epsilon} --- Spacing near one supplies a reference precision scale, not a universal local gap or tolerance.
\item
  \hyperref[appendix-14bab4cab9d4]{ULP and representable spacing} --- Adjacent representable points define rounding candidates, errors and step distances.
\item
  \hyperref[appendix-eb336149c00a]{Floating-point infinities} --- Non-finite operands and overflowing intermediates need branches before ordinary finite formulas.
\item
  \hyperref[appendix-d6607cc4c235]{NaNs and payloads} --- Unordered comparisons and payload distinctions require explicit metric, comparison and serialization policy.
\end{itemize}

\section{Used by}\label{appendix-0ec599f8ba27-used-by}

\begin{itemize}
\item
  \href{../02-Foundations-II/main.pdf\#nameddest=appendix-b5b83329215b}{Vector tolerances and finite-state diagnostics (02-Foundations-II, appendix-b5b83329215b)} --- Phase 6 uses this established foundation in its geometry narrative.
\item
  \hyperref[appendix-8f6a01f284a7]{Conditioning and singularity} --- uses this mechanism in its explanation.
\end{itemize}

\section{Recall and next depth}\label{appendix-0ec599f8ba27-recall-and-next-depth}

Reconstruct the concrete example in the definition and identify its validity assumptions.

\textbf{Pending:} Extend this mechanism to later algorithms and target-specific behavior. See the M005 chapter for its current examples and validity assumptions.

\section{Source and navigation}\label{appendix-0ec599f8ba27-source-and-navigation}

\href{Continuous\%20Notes/Computing\%20Foundations\%20-\%20From\%20Bits\%20to\%20C++\%20and\%20Systems\%20APIs\#When\%20arithmetic\%20meets\%20uncertainty}{M005 development}

\href{Continuous\%20Notes/Computing\%20Foundations\%20-\%20From\%20Bits\%20to\%20C++\%20and\%20Systems\%20APIs\#A\%20stretching\%20ruler\%20for\%20real-valued\%20quantities}{M004 teaching chapter} · \hyperref[appendix-fd090d982a15-m004---parts-61-76]{Per-part definitions} · \href{../02-Foundations-II/main.pdf\#nameddest=appendix-4ad47bdd269a}{Dependency Map (02-Foundations-II, appendix-4ad47bdd269a)}
